\chapter{HASIL DAN PEMBAHASAN}
\label{ch:hasil}

Bab ini melaporkan hasil eksperimen dan implementasi yang dirancang pada Bab~\ref{ch:metode}. Pemaparan diawali dengan lingkungan pengujian, dilanjutkan dengan persiapan dan eksplorasi dataset, hasil eksperimen utama menggunakan kombinasi DINOv3 \textit{large}, \textit{multilingual-E5 large}, dan CatBoost yang dikalibrasi PGS, hasil ablasi modalitas, hasil ablasi kalibrasi, hasil ablasi pasangan \textit{encoder}, validasi paritas ONNX, hasil implementasi aplikasi Flutter beserta integrasi Supabase, kemudian diakhiri dengan pembahasan dan analisis kesalahan.

% =====================================================================
\section{Lingkungan Pengujian}
\label{sec:lingkungan-pengujian}
% =====================================================================

Eksperimen pelatihan dan ekstraksi \textit{embedding} dijalankan pada lingkungan komputasi awan RunPod yang diakses melalui SSH, dengan akselerator GPU yang dipilih sesuai ketersediaan. Lingkungan perangkat lunak yang dipakai dirangkum pada Tabel~\ref{tab:lingkungan-software}.

\begin{table}[htbp]
\centering
\caption{Lingkungan perangkat lunak utama.}
\label{tab:lingkungan-software}
\begin{tabular}{p{4.5cm} p{8.0cm}}
\toprule
\textbf{Komponen} & \textbf{Versi atau Detail} \\
\midrule
Python & 3.11 \\
PyTorch & 2.x dengan CUDA toolkit \\
ONNX dan ONNX Runtime & 1.x untuk ekspor dan inferensi \\
CatBoost & 1.2.x dengan opsi \textit{posterior sampling} \\
FastAPI & 0.x untuk server inferensi \\
Flutter & 3.x untuk klien Android \\
Supabase Dart Client & versi stabil terbaru \\
RunPod & GPU \textit{cloud} NVIDIA GeForce RTX 4090 (24 GB VRAM), sesuai Tabel~\ref{tab:runpod} \\
\bottomrule
\end{tabular}
\end{table}

Pengujian aplikasi Flutter dilakukan pada perangkat Android fisik dengan koneksi USB \textit{tethering}, dengan port server FastAPI yang diteruskan melalui \texttt{adb reverse} sehingga \textit{endpoint} \texttt{/predict} dapat diakses oleh aplikasi seolah-olah berjalan di perangkat lokal. Pada \textit{deployment} publik, server yang sama dijalankan sebagai HuggingFace Space berbasis Docker dengan URL HTTPS otomatis sehingga dapat dijangkau tanpa \textit{tethering}.

% =====================================================================
\section{Persiapan dan Eksplorasi Data}
\label{sec:persiapan-data}
% =====================================================================

Setelah pemetaan 487 nama SKPD mentah ke sembilan kategori target dan pembagian stratifikasi, partisi akhir dataset berisi 43.241 sampel pelatihan, 9.266 sampel validasi, dan 9.266 sampel uji. Ketidakseimbangan kelas tidak dihilangkan melalui \textit{resampling}, melainkan dikompensasi melalui pembobotan kelas (\textit{class weights}) pada algoritma CatBoost. Distribusi sampel pelatihan per instansi dirangkum pada Gambar~\ref{fig:class-distribution} dan Tabel~\ref{tab:class-distribution}.

\begin{figure}[htbp]
    \centering
    \includegraphics[width=0.92\textwidth]{bab_4/class_distribution.png}
    \caption{Distribusi sembilan instansi (dinas) pada set train.}
    \label{fig:class-distribution}
\end{figure}

\begin{table}[htbp]
\centering
\caption{Distribusi sampel per instansi (dinas) pada set train.}
\label{tab:class-distribution}
\small\setlength{\tabcolsep}{4pt}
\begin{tabular}{c l r r}
\toprule
\textbf{ID} & \textbf{Kategori} & \textbf{Jumlah Sampel} & \textbf{Proporsi (\%)} \\
\midrule
0 & Dinas Bina Marga & 9.467 & 21,89 \\
1 & Satuan Polisi Pamong Praja & 5.883 & 13,60 \\
2 & Dinas Perhubungan & 8.175 & 18,90 \\
3 & Kelurahan & 8.386 & 19,39 \\
4 & Dinas Pertamanan dan Hutan & 4.200 & 9,71 \\
5 & Dinas Sumber Daya Air & 1.823 & 4,22 \\
6 & Dinas Cipta Karya, Tata Ruang, dan Pertanahan & 1.033 & 2,39 \\
7 & Badan Pembinaan BUMD & 1.179 & 2,73 \\
8 & Instansi lain & 3.095 & 7,16 \\
\midrule
\multicolumn{2}{l}{\textbf{Total}} & \textbf{43.241} & \textbf{100,00} \\
\bottomrule
\end{tabular}
\end{table}

Distribusi yang ditampilkan adalah distribusi apa adanya pada set \textit{train} tanpa \textit{resampling}; ketidakseimbangan sengaja tidak dihapus dari data, melainkan dikompensasi melalui pembobotan kelas saat pelatihan sebagaimana diuraikan pada Bagian~\ref{sec:bobot-kelas}. Empat kelas teratas (Dinas Bina Marga, Kelurahan, Dinas Perhubungan, dan Satpol PP) mencakup lebih dari 73 persen total sampel, sedangkan tiga kelas terkecil (Dinas Cipta Karya, Tata Ruang, dan Pertanahan, Badan Pembinaan BUMD, serta Dinas Sumber Daya Air) hanya mencakup sekitar 9 persen. Ketidakseimbangan ini menjadi pertimbangan saat pemilihan metrik utama, yaitu \textit{macro F1 score} yang memberikan bobot sama pada setiap kelas, sebagaimana diuraikan pada Bagian~\ref{sec:skenario-pengujian}.

% =====================================================================
\section{Hasil Eksperimen Utama}
\label{sec:hasil-utama}
% =====================================================================

Konfigurasi terbaik yang dipilih untuk \textit{deployment} adalah kombinasi \textit{image encoder} DINOv3 \textit{large}, \textit{text encoder multilingual-E5 large}, fusi awal berbasis konkatenasi, \textit{classifier head} CatBoost, dan kalibrasi probabilitas dengan PGS. Tabel~\ref{tab:hasil-utama} merangkum metrik kinerja konfigurasi tersebut pada set uji.

\begin{table}[htbp]
\centering
\caption{Metrik kinerja konfigurasi terbaik pada set uji (DINOv3 large + mE5 large + CatBoost+PGS).}
\label{tab:hasil-utama}
\begin{tabular}{l c c}
\toprule
\textbf{Konfigurasi} & \textbf{Akurasi} & \textbf{Macro F1} \\
\midrule
CatBoost (tanpa PGS) & 0,7996 & 0,7684 \\
CatBoost dengan PGS & \textbf{0,8074} & \textbf{0,7747} \\
$\Delta$ (PGS minus tanpa PGS) & +0,0078 & +0,0063 \\
\bottomrule
\end{tabular}
\end{table}

Konfigurasi terbaik mencapai \textit{akurasi} 0,8074 dan \textit{macro F1} 0,7747 pada set uji 9.266 sampel, dengan \textit{balanced accuracy} (rata-rata \textit{recall} per kelas) 0,7916 dan \textit{macro ROC-AUC} (\textit{one-vs-rest}) sebesar 0,9667. Capaian ini kompetitif terhadap baseline yang dilaporkan oleh \citet{hanifPenugasanOtomatis2024} sebesar akurasi 80,5 persen pada delapan kelas. Perbandingan ini perlu dibaca secara hati-hati dan bukan sebagai unggul-langsung pada kondisi eksperimen yang identik, karena jumlah kelas berbeda: penelitian ini melaporkan \textit{akurasi}, \textit{balanced accuracy}, dan \textit{macro F1} pada sembilan kelas (termasuk kelas \textit{catch-all} Instansi lain yang lebih sulit), sedangkan baseline melaporkan akurasi pada delapan kelas. Dengan kondisi tugas yang lebih luas tersebut, penggantian DINOv2 dengan DINOv3 \textit{large} tetap mempertahankan kinerja pada tingkat yang setara atau sedikit lebih tinggi. Rerata waktu pelatihan satu konfigurasi CatBoost adalah 308 detik tanpa PGS dan 602 detik dengan PGS, yang membuktikan bahwa biaya komputasi pelatihan masih dapat ditoleransi pada satu sesi eksperimen.

\subsection{Metrik Per-Kelas}
\label{sec:metrik-per-kelas}

Untuk mengamati kinerja model pada masing-masing instansi, Tabel~\ref{tab:per-class} menyajikan \textit{precision}, \textit{recall}, dan \textit{F1-score} per kelas pada set uji untuk konfigurasi terbaik (DINOv3 \textit{large} + mE5 \textit{large} + CatBoost+PGS).

\begin{table}[htbp]
\centering
\caption{Metrik per-kelas pada set uji (DINOv3 large + mE5 large + CatBoost+PGS).}
\label{tab:per-class}
\small\setlength{\tabcolsep}{4pt}
\begin{tabular}{l c c c r}
\toprule
\textbf{Kategori (Dinas)} & \textbf{Precision} & \textbf{Recall} & \textbf{F1} & \textbf{Support} \\
\midrule
Dinas Bina Marga & 0,8644 & 0,8797 & 0,8720 & 2.029 \\
Satuan Polisi Pamong Praja & 0,8432 & 0,8699 & 0,8564 & 1.261 \\
Dinas Perhubungan & 0,8443 & 0,8699 & 0,8569 & 1.752 \\
Kelurahan & 0,8103 & 0,6939 & 0,7476 & 1.797 \\
Dinas Pertamanan dan Hutan & 0,8924 & 0,9122 & \textbf{0,9022} & 900 \\
Dinas Sumber Daya Air & 0,5577 & 0,7417 & 0,6367 & 391 \\
Dinas Cipta Karya, Tata Ruang, dan Pertanahan & 0,8398 & 0,8778 & 0,8584 & 221 \\
Badan Pembinaan BUMD & 0,6711 & 0,7937 & 0,7273 & 252 \\
Instansi lain & 0,5486 & 0,4857 & \textbf{0,5152} & 663 \\
\midrule
\textbf{Macro average} & 0,7635 & 0,7916 & \textbf{0,7747} & 9.266 \\
\textbf{Weighted average} & 0,8086 & 0,8073 & 0,8061 & 9.266 \\
\bottomrule
\end{tabular}
\end{table}

Tabel per-kelas mengungkap bahwa kelas dengan dukungan besar dan kata kunci spesifik (Dinas Pertamanan dan Hutan, $F1 = 0{,}9022$; Dinas Bina Marga, $F1 = 0{,}8720$) terklasifikasi paling baik. Sebaliknya, kelas \textit{catch-all} ``Instansi lain'' ($F1 = 0{,}5152$) menjadi yang terlemah karena memuat laporan heterogen tanpa pola visual atau tekstual yang konsisten, diikuti Dinas Sumber Daya Air ($F1 = 0{,}6367$) yang memiliki \textit{precision} rendah (0,5577) akibat tertukar dengan kelas lain yang juga berkaitan dengan saluran air dan banjir. Menariknya, dua kelas minoritas lain yaitu Dinas Cipta Karya ($F1 = 0{,}8584$, hanya 221 sampel) dan Badan Pembinaan BUMD ($F1 = 0{,}7273$, 252 sampel) tetap terklasifikasi cukup baik, yang menunjukkan bahwa kelangkaan sampel tidak selalu berarti kinerja buruk apabila pola laporannya khas.

% =====================================================================
\section{Hasil Ablasi Modalitas}
\label{sec:ablasi-modalitas}
% =====================================================================

Untuk memvalidasi kontribusi setiap modalitas, eksperimen ablasi dijalankan dengan tiga konfigurasi, yaitu citra saja, teks saja, dan fusi awal. Hasilnya dirangkum pada Tabel~\ref{tab:ablasi-modalitas} dan Gambar~\ref{fig:modality-ablation}.

\begin{table}[htbp]
\centering
\caption{Hasil ablasi modalitas pada set uji.}
\label{tab:ablasi-modalitas}
\begin{tabular}{l c c c c}
\toprule
\textbf{Konfigurasi} & \textbf{Acc CatBoost} & \textbf{F1 CatBoost} & \textbf{Acc PGS} & \textbf{F1 PGS} \\
\midrule
Citra saja (DINOv3 large) & 0,6265 & 0,6016 & 0,6238 & 0,5973 \\
Teks saja (mE5 large) & 0,7850 & 0,7498 & 0,7874 & 0,7515 \\
Fusi awal (citra + teks) & \textbf{0,7996} & \textbf{0,7684} & \textbf{0,8074} & \textbf{0,7747} \\
\bottomrule
\end{tabular}
\end{table}

\begin{figure}[htbp]
    \centering
    \includegraphics[width=0.92\textwidth]{bab_4/modality_ablation.png}
    \caption{Ablasi modalitas pada CatBoost dan CatBoost+PGS.}
    \label{fig:modality-ablation}
\end{figure}

Hasil ini menyampaikan tiga temuan penting. Pertama, modalitas teks saja jauh lebih diskriminatif dibandingkan modalitas citra saja, yaitu \textit{macro F1} 0,7498 untuk teks dibandingkan 0,6016 untuk citra. Hal ini wajar karena narasi laporan sering kali memuat petunjuk eksplisit tentang dinas yang dituju, sedangkan citra hanya menampilkan kondisi visual yang ambigu pada beberapa kategori. Kedua, fusi awal memberikan peningkatan yang konsisten di atas modalitas tunggal terbaik, yaitu peningkatan \textit{macro F1} sebesar 0,0232 (2,32 \textit{percentage point}) terhadap teks saja pada konfigurasi PGS, menunjukkan bahwa citra membawa informasi tambahan yang tidak selalu tertangkap oleh teks. Ketiga, pengaruh PGS pada modalitas citra saja adalah negatif tipis (selisih $-0,0043$ pada F1), sedangkan pada modalitas teks saja dan fusi awal pengaruhnya positif. Hal ini mengindikasikan bahwa PGS bekerja optimal ketika model awal sudah memberikan probabilitas yang relatif terkalibrasi, kondisi yang lebih sulit dipenuhi oleh model yang hanya mengandalkan modalitas visual.

% =====================================================================
\section{Hasil Ablasi Kalibrasi (PGS)}
\label{sec:ablasi-pgs}
% =====================================================================

Pengaruh kalibrasi probabilitas dengan PGS dilaporkan pada Tabel~\ref{tab:ablasi-pgs} dan divisualisasikan pada Gambar~\ref{fig:pgs-comparison}. Konfigurasi yang dievaluasi adalah pasangan \textit{encoder} terbaik (DINOv3 \textit{large} + mE5 \textit{large}) dengan dan tanpa PGS.

\begin{table}[htbp]
\centering
\caption{Hasil ablasi kalibrasi pada set uji.}
\label{tab:ablasi-pgs}
\begin{tabular}{l c c}
\toprule
\textbf{Model} & \textbf{Akurasi} & \textbf{Macro F1} \\
\midrule
CatBoost (tanpa PGS) & 0,7996 & 0,7684 \\
CatBoost dengan PGS & \textbf{0,8074} & \textbf{0,7747} \\
$\Delta$ (PGS minus tanpa PGS) & +0,0078 & +0,0063 \\
\bottomrule
\end{tabular}
\end{table}

\begin{figure}[htbp]
    \centering
    \includegraphics[width=0.78\textwidth]{bab_4/pgs_comparison.png}
    \caption{Perbandingan akurasi dan macro F1 dengan dan tanpa PGS.}
    \label{fig:pgs-comparison}
\end{figure}

PGS memberikan peningkatan \textit{akurasi} sebesar 0,78 \textit{percentage point} dan \textit{macro F1} sebesar 0,63 \textit{percentage point} terhadap baseline tanpa PGS. Selain memberikan peningkatan akurasi, PGS juga menghasilkan ukuran ketidakpastian (\textit{mean uncertainty}) sebesar 0,00339 yang dapat ditampilkan pada antarmuka aplikasi sebagai sinyal kepercayaan kepada warga pelapor. Sinyal ini menjadi penting saat warga harus memutuskan apakah hasil klasifikasi otomatis perlu dikoreksi secara manual sebelum laporan dikirim.

\subsection{Mengapa Peningkatan PGS Tergolong Marginal}
\label{sec:pgs-marginal}

Peningkatan dari PGS tergolong kecil ($+0{,}63$ pp \textit{macro F1}) dan hal ini konsisten dengan sifat metode serta karakteristik data, bukan indikasi kegagalan implementasi. Empat penjelasan dapat dikemukakan.

Pertama, PGS pada hakikatnya merupakan estimator \textit{posterior mean} di atas ensembel virtual CatBoost \citep{ustimenkoCatBoostUncertainty2023}; nilai harapan posterior tersebut sangat dekat dengan prediksi titik model tunggal, sehingga label \textit{argmax} hanya berubah pada sampel yang berada tepat di sekitar batas keputusan. Karena mayoritas sampel uji diklasifikasikan dengan \textit{margin} lebar, tercermin dari \textit{macro ROC-AUC} 0,9667, jumlah sampel yang berpindah label akibat PGS sedikit, sehingga selisih akurasi pun kecil.

Kedua, sumber ketidakpastian dominan pada dataset ini bersifat \textit{aleatoric} (melekat pada data), bukan \textit{epistemic} (akibat keterbatasan model). PGS dirancang untuk merefinasi komponen \textit{epistemic} melalui rerata ensembel, sedangkan ambiguitas dominan di sini berasal dari pemetaan \textit{many-to-one} 487 nama SKPD ke sembilan kelas, narasi laporan yang sangat singkat, dan laporan yang secara historis sah dirutekan ke lebih dari satu instansi. Ketidakpastian \textit{aleatoric} bersifat irreducibel sehingga tidak dapat dikurangi oleh penambahan anggota ensembel \citep{kendallGal2017,huellermeierWaegeman2021}. Hal ini sekaligus menjelaskan mengapa \textit{mean epistemic uncertainty} yang terukur sangat kecil, yaitu 0,00339.

Ketiga, ensembel virtual PGS berbagi parameter dari satu model CatBoost yang sama, bukan model independen seperti pada \textit{deep ensembles} \citep{lakshminarayananDeepEnsembles2017}. Korelasi antar-anggota yang tinggi membuat sebaran \textit{epistemic} yang dilaporkan cenderung bias rendah, sehingga ruang perbaikan dari agregasi posterior memang sempit. Konsisten dengan ini, ablasi modalitas menunjukkan PGS justru sedikit menurunkan F1 pada modalitas citra-saja ($-0{,}43$ pp), yaitu kondisi ketika probabilitas dasar kurang terkalibrasi sehingga rerata ensembel memperkuat bias alih-alih memperbaikinya.

Pola yang sama dapat diamati langsung di dalam eksperimen penelitian ini sendiri, tanpa perlu bersandar pada perbandingan lintas-penelitian. Pada matrix grid pasangan \textit{encoder} (Tabel~\ref{tab:ablasi-encoder}), \textit{macro F1} setelah kalibrasi PGS nyaris datar di antara \textit{image encoder} yang sudah kompeten, yaitu 0,7747 untuk DINOv3 \textit{large}, 0,7747 untuk EVA-02 \textit{large}, dan 0,7736 untuk DINOv2 \textit{large}. Keseragaman ini menunjukkan bahwa begitu representasi dasar mencapai tingkat keterpisahan kelas yang memadai, baik pemilihan \textit{encoder} citra maupun agregasi posterior PGS tidak lagi banyak menggeser kinerja; batas atas justru ditentukan oleh faktor \textit{aleatoric} pada data, bukan oleh penghalusan \textit{epistemic}. Hal ini secara struktural memang tak terhindarkan: karena \textit{posterior mean} PGS hanya bergeser tipis dari prediksi titik model tunggal, besar perbaikannya dibatasi oleh banyaknya sampel yang berada tepat di sekitar batas keputusan, dan jumlah itu kecil ketika \textit{margin} klasifikasi lebar (tercermin dari \textit{macro ROC-AUC} 0,9667). Dengan kata lain, semakin baik dan semakin terkalibrasi representasi dasar, semakin kecil manfaat marjinal dari kalibrasi ensembel virtual, persis seperti yang diprediksi oleh argumen \textit{aleatoric-epistemic} di atas.

Dengan demikian, kontribusi utama PGS pada penelitian ini bukan terletak pada lonjakan akurasi, melainkan pada penyediaan sinyal ketidakpastian berbiaya inferensi nyaris nol yang dapat ditampilkan pada antarmuka sebagai dasar koreksi manual. Untuk memperoleh keandalan \textit{end-to-end} yang lebih kuat, arah lanjutan yang lebih tepat adalah lapisan \textit{selective prediction} berjaminan cakupan atau keluaran himpunan-bernilai melalui \textit{conformal prediction}, ketimbang sekadar menambah anggota ensembel virtual.

Penilaian signifikansi statistik memperkuat kesimpulan ini. Dengan \textit{bootstrap} sebanyak 2.000 kali pengambilan ulang (\textit{resampling}) pada 9.266 sampel uji, \textit{macro F1} konfigurasi CatBoost+PGS sebesar 0,7747 memiliki selang kepercayaan 95\% sebesar $[0{,}764;\ 0{,}784]$. Karena \textit{macro F1} baseline tanpa PGS (0,7684) berada di dalam selang tersebut, peningkatan $+0{,}63$ pp tidak dapat diklaim signifikan secara statistik pada taraf 95\%. Temuan ini konsisten dengan posisi bahwa nilai tambah PGS terletak pada penyediaan estimasi ketidakpastian, bukan pada peningkatan akurasi yang dapat dibedakan dari ragam pengambilan sampel.

% =====================================================================
\section{Hasil Ablasi Pasangan Encoder}
\label{sec:ablasi-encoder}
% =====================================================================

Tabel~\ref{tab:ablasi-encoder} dan Gambar~\ref{fig:encoder-matrix} merangkum hasil ablasi pasangan \textit{image encoder} dan \textit{text encoder} dengan kalibrasi PGS aktif. Total enam belas pasangan dievaluasi (empat \textit{image encoder} dikali empat \textit{text encoder}).

\begin{table}[htbp]
\centering
\small
\caption{Macro F1 (CatBoost+PGS) pada matrix grid pasangan encoder.}
\label{tab:ablasi-encoder}
\begin{tabular}{l c c c c}
\toprule
\textbf{Image $\backslash$ Text} & \textbf{mE5 large} & \textbf{BGE-M3} & \textbf{IndoBERT} & \textbf{Cendol-mT5} \\
\midrule
DINOv3 large  & \textbf{0,7747} & 0,7628 & 0,7449 & 0,7340 \\
DINOv2 large  & 0,7736 & 0,7651 & 0,7471 & 0,7305 \\
EVA-02 large  & 0,7747 & 0,7634 & 0,7429 & 0,7323 \\
Hiera large   & 0,7652 & 0,7558 & 0,7285 & 0,7149 \\
\bottomrule
\end{tabular}
\end{table}

\begin{figure}[htbp]
    \centering
    \includegraphics[width=0.86\textwidth]{bab_4/encoder_matrix.png}
    \caption{Heatmap macro F1 pada matrix grid pasangan encoder.}
    \label{fig:encoder-matrix}
\end{figure}

Pengamatan utama dari matrix grid ini adalah sebagai berikut. Pertama, \textit{text encoder multilingual-E5 large} secara konsisten menjadi pilihan terbaik di seluruh empat \textit{image encoder}, mengungguli BGE-M3, IndoBERT, dan Cendol-mT5 pada setiap baris. Hal ini sejalan dengan kekuatan multilingual-E5 dalam menyandikan kalimat Bahasa Indonesia yang informal dan singkat ke ruang vektor yang diskriminatif. Kedua, di antara empat \textit{image encoder}, DINOv3 \textit{large} dan EVA-02 \textit{large} menunjukkan kinerja yang nyaris identik pada \textit{macro F1} (0,7747), keduanya unggul terhadap DINOv2 \textit{large} dan Hiera \textit{large}. Penelitian ini memilih DINOv3 \textit{large} untuk \textit{deployment} karena merupakan keluarga \textit{encoder} paling mutakhir pada saat eksperimen. Ketiga, Hiera \textit{large} secara konsisten menjadi yang terlemah pada keempat pasangan, yang menunjukkan bahwa arsitektur hierarchical yang dioptimalkan untuk efisiensi inferensi belum mengejar kualitas representasi DINOv3 atau EVA-02 pada domain laporan publik perkotaan.

% =====================================================================
\section{Validasi Paritas ONNX}
\label{sec:paritas-onnx}
% =====================================================================

Setelah \textit{checkpoint} CatBoost terbaik dipilih, model \textit{native} \texttt{.cbm} diekspor ke format ONNX bersama dengan \textit{image encoder} DINOv3 \textit{large} dan \textit{text encoder multilingual-E5 large}. Validasi paritas dilakukan dengan menjalankan kedua jalur (PyTorch dan ONNX) pada batch sampel uji yang sama, lalu menghitung selisih probabilitas absolut maksimum di seluruh komponen vektor keluaran. Hasilnya dirangkum pada Tabel~\ref{tab:paritas-onnx}.

\begin{table}[htbp]
\centering
\caption{Validasi paritas numerik PyTorch terhadap ONNX pada set uji.}
\label{tab:paritas-onnx}
\begin{tabular}{l c}
\toprule
\textbf{Komponen} & \textbf{$|\Delta p|_{\max}$} \\
\midrule
\textit{Image encoder} (DINOv3 large) & $< 10^{-5}$ \\
\textit{Text encoder} (multilingual-E5 large) & $< 10^{-5}$ \\
\textit{Classifier head} (CatBoost+PGS) & $< 10^{-5}$ \\
End-to-end (gabungan ketiga komponen) & $< 10^{-5}$ \\
\bottomrule
\end{tabular}
\end{table}

Selisih probabilitas absolut maksimum di bawah ambang $10^{-5}$ pada seluruh komponen memastikan prediksi top-1 identik antara jalur PyTorch dan jalur ONNX, sehingga server FastAPI yang memuat ketiga \textit{artifact} ONNX dapat dipakai sebagai jalur produksi tanpa kekhawatiran terjadi divergensi numerik dengan jalur pelatihan.

% =====================================================================
\section{Hasil Implementasi Aplikasi Flutter}
\label{sec:hasil-flutter}
% =====================================================================

Aplikasi SmartCityApps berhasil diimplementasikan menggunakan Flutter dengan arsitektur berbasis fitur (\textit{feature-first}), manajemen \textit{state} Riverpod, navigasi deklaratif GoRouter, dan \textit{backend} Supabase (PostgreSQL, Auth, Storage, dan Realtime). Inferensi klasifikasi tidak dijalankan di perangkat melainkan melalui \textit{endpoint} \texttt{/predict} pada server FastAPI, sehingga aplikasi tetap ringan dan model dapat diperbarui terpusat tanpa pembaruan aplikasi. Aplikasi mendukung tiga peran pengguna, yaitu Warga Pelapor, Operator Instansi, dan Moderator atau \textit{Super Admin}. Pemilahan peran ditegakkan pada dua lapis: lapis klien memilih pohon navigasi dan ruang-lingkup data (\textit{citizen}, \textit{agency}, atau \textit{global scope}) berdasarkan peran, sedangkan lapis basis data menegakkan otorisasi sesungguhnya melalui \textit{Row Level Security} (RLS) pada tabel \texttt{profiles}, \texttt{reports}, \texttt{report\_history}, dan \texttt{agencies}. Otentikasi memanfaatkan Supabase Auth dengan konfirmasi surel dan pemulihan kata sandi melalui \textit{deep link}, sementara pembaruan status laporan dipropagasikan ke antarmuka secara langsung melalui Postgres Realtime tanpa \textit{polling}.

\subsection{Tangkapan Layar Antarmuka Pengguna}
\label{sec:ui-screenshots}

Alur inti pelaporan oleh Warga Pelapor ditampilkan pada Gambar~\ref{fig:ui-pelaporan}, dan tampilan peran Moderator atau \textit{Super Admin} pada Gambar~\ref{fig:ui-superadmin}; galeri lengkap seluruh layar untuk ketiga peran (onboarding, autentikasi, beranda, peta, pengaturan, detail laporan, dasbor instansi, hingga registri instansi) disajikan pada Lampiran~\ref{app:ui-mockup}. Tangkapan layar bersumber dari dua kanal yang saling melengkapi. Alur peran Warga Pelapor ditangkap langsung dari perangkat Android fisik saat aplikasi berjalan \textit{end-to-end} terhadap server inferensi dan Supabase, sehingga memuat data nyata (foto kejadian, alamat hasil \textit{reverse-geocoding}, nilai \textit{confidence}, instansi tujuan beserta estimasi jarak, dan metadata model). Layar yang aksesnya dibatasi peran Operator Instansi dan Moderator dirender secara deterministik melalui \textit{golden screenshot test} (pengujian render \textit{widget} Flutter yang sebenarnya dengan data contoh terkendali), karena akun berperan instansi belum di-\textit{provisioning} pada perangkat uji; pendekatan ini menjaga tampilan tetap konsisten dan dapat direproduksi tanpa bergantung pada koneksi jaringan.

\begin{figure}[htbp]
    \centering
    \begin{subfigure}[b]{0.31\textwidth}
        \centering
        \includegraphics[height=6.4cm]{image-androdid/buat-laporan.jpg}
        \caption{Buat laporan}
        \label{fig:ui-create}
    \end{subfigure}
    \hfill
    \begin{subfigure}[b]{0.31\textwidth}
        \centering
        \includegraphics[height=6.4cm]{image-androdid/Hasil-analisis-ai.jpg}
        \caption{Tinjauan hasil AI}
        \label{fig:ui-aireview}
    \end{subfigure}
    \hfill
    \begin{subfigure}[b]{0.31\textwidth}
        \centering
        \includegraphics[height=6.4cm]{image-androdid/melihat-laporan.jpg}
        \caption{Riwayat laporan}
        \label{fig:ui-myreports}
    \end{subfigure}
    \caption{Alur inti pelaporan pada perangkat Android fisik: pengambilan foto, narasi, dan akuisisi lokasi GPS (a); layar tinjauan AI yang menampilkan instansi prediksi beserta \textit{confidence} ($88\%$), peringkat prediksi alternatif, instansi tujuan terdekat dengan estimasi jarak, dan opsi koreksi kategori manual (b); serta daftar riwayat laporan beserta status terkini (c). Galeri lengkap seluruh layar disajikan pada Lampiran~\ref{app:ui-mockup}.}
    \label{fig:ui-pelaporan}
\end{figure}

\subsection{Skema Basis Data Supabase}

Aplikasi ditopang skema ternormalisasi dengan \texttt{reports} sebagai akar agregat (laporan dengan instansi prediksi, confidence, probabilitas multiclass, lokasi, status, visibilitas, dan penanda duplikat) yang dilengkapi tabel anak \texttt{report\_evidence} (bukti foto, banyak per laporan), \texttt{report\_assignments} (riwayat penugasan dan penugasan-ulang instansi), \texttt{report\_history} (audit status), dan \texttt{report\_sla} (tenggat layanan); tabel \texttt{agencies} menjadi registri instansi tujuan dan \texttt{profiles} menyimpan data pengguna serta peran. Berkas disimpan pada \textit{bucket} privat Supabase Storage sebagai pasangan \textit{bucket} dan \texttt{object\_path}, lalu diakses melalui \textit{signed URL} berumur singkat. Agar klien Flutter tetap membaca satu tabel melalui Supabase Postgres Realtime, kolom datar pada \texttt{reports} dipertahankan sebagai \textit{proyeksi-baca} yang disinkronkan \textit{trigger} dari tabel anak, sehingga perubahan status langsung tercermin di antarmuka tanpa \textit{polling}. Skema yang benar-benar ter-deploy pada Supabase ditunjukkan pada Gambar~\ref{fig:erd-supabase}, sedangkan rancangan konseptualnya pada Gambar~\ref{fig:db-schema}.

\begin{sidewaysfigure}
    \centering
    \includegraphics[width=0.95\textheight]{supabase/public_schema_deployed.png}
    \caption{Skema basis data yang ter-deploy pada Supabase (hasil visualisasi Supabase Studio): tabel publik aplikasi setelah normalisasi, yaitu akar agregat \texttt{reports} dengan tabel anak \texttt{report\_evidence}, \texttt{report\_assignments}, \texttt{report\_history}, \texttt{report\_sla}, dan registri \texttt{agencies}. Rancangan konseptualnya beserta mekanisme proyeksi-baca disajikan pada Gambar~\ref{fig:db-schema}.}
    \label{fig:erd-supabase}
\end{sidewaysfigure}

\subsection{Alur Pelaporan oleh Warga}

Skenario alur pelaporan oleh Warga Pelapor diuji secara end-to-end pada perangkat Android fisik. Tahapan yang berhasil dijalankan adalah pengambilan foto melalui kamera atau galeri, pengisian narasi laporan, akuisisi koordinat GPS dengan opsi pemilihan manual di peta bila layanan lokasi tidak tersedia, klasifikasi otomatis melalui \textit{endpoint} \texttt{/predict} server FastAPI, tinjauan hasil pada layar tinjauan AI, opsi koreksi kategori secara manual, tampilan instansi tujuan terdekat dari backend, dan terakhir penyimpanan laporan ke tabel \texttt{reports} di Supabase beserta unggahan gambar ke Storage.

\subsection{Alur Verifikasi oleh Operator Instansi}

Operator Instansi diarahkan ke dasbor admin (antrean \texttt{AgencyScope}) yang, melalui RLS, hanya menampilkan laporan ter-\textit{routing} ke instansi yang ditugaskan kepada operator tersebut. Antrean dilengkapi panel \textit{triage} yang menonjolkan jumlah laporan berstatus \textit{submitted} (menunggu tindakan) dan bilah tapis status (\textit{submitted}, \textit{verified}, \textit{in progress}, \textit{resolved}) sehingga operator dapat memprioritaskan beban kerja. Saat membuka satu laporan, operator memperoleh bilah aksi pembaruan status: memverifikasi laporan, menolak dengan mengisi alasan penolakan, menandai sedang diproses, atau menandai selesai. Transisi ke status selesai diwajibkan menyertakan unggah foto bukti penyelesaian sebagai prasyarat (\textit{gating}), sehingga klaim penyelesaian selalu terdokumentasi. Setiap perubahan status menulis baris audit ke tabel \texttt{report\_history}, dan kolom datar pada \texttt{reports} disinkronkan melalui \textit{trigger} sehingga pelapor menerima pembaruan status secara \textit{realtime}.

\subsection{Alur Moderasi dan Registri Instansi}

Moderator atau \textit{Super Admin} memiliki hak akses lintas instansi (\texttt{GlobalScope}). Pada antrean global ``All agencies'' (Gambar~\ref{fig:ui-superadmin}a), moderator dapat meninjau seluruh laporan dari semua instansi sekaligus, menyaringnya per instansi melalui tapis tambahan yang khusus muncul bagi peran ini, serta memantau panel \textit{triage} lintas-instansi. Selain moderasi laporan, peran ini mengelola registri instansi tujuan (Gambar~\ref{fig:ui-superadmin}b), yaitu antarmuka CRUD atas tabel \texttt{agencies} yang menjadi \textit{system of record} bagi algoritma perutean instansi terdekat. Setiap kantor memuat nama, kategori, koordinat, dan penanda aktif; aksi ``Terapkan'' memanggil \textit{endpoint} \texttt{/reload-agencies} pada server inferensi agar perutean menggunakan data terbaru tanpa \textit{restart} penuh. Penulisan ke registri dibatasi peran \texttt{super\_admin} melalui kebijakan RLS \texttt{agencies\_write\_super\_admin}. Pengujian alur moderasi memastikan RLS berfungsi dengan benar, yaitu Moderator dapat mengakses tabel lintas instansi sedangkan Operator hanya dapat mengakses laporan pada instansi yang ditugaskan, dan Warga hanya dapat mengakses laporannya sendiri.

\begin{figure}[htbp]
    \centering
    \begin{subfigure}[b]{0.31\textwidth}
        \centering
        \includegraphics[width=\linewidth]{bab_3/ui/ui_superadmin_inbox.png}
        \caption{Antrean lintas instansi}
        \label{fig:ui-superadmin-inbox}
    \end{subfigure}
    \hspace{0.06\textwidth}
    \begin{subfigure}[b]{0.31\textwidth}
        \centering
        \includegraphics[width=\linewidth]{bab_3/ui/ui_agency_registry.png}
        \caption{Registri instansi}
        \label{fig:ui-superadmin-registry}
    \end{subfigure}
    \caption{Tampilan peran Moderator atau \textit{Super Admin}: antrean global ``All agencies'' yang menggabungkan laporan dari seluruh instansi dengan tapis instansi dan panel \textit{triage} (a); serta layar registri instansi (CRUD atas tabel \texttt{agencies}) yang memuat nama kantor, kategori, koordinat, penanda aktif, dan aksi pemuatan ulang server perutean (b).}
    \label{fig:ui-superadmin}
\end{figure}

\subsection{Demo Perutean Instansi Terdekat}

Mesin perutean memilih instansi tujuan terdekat berdasarkan instansi prediksi dan koordinat laporan. Tabel~\ref{tab:demo-perutean} menampilkan tiga contoh laporan beserta instansi prediksi dan instansi tujuan hasil perutean. Entri pada kolom tujuan bersifat ilustratif pada tingkat kategori; pemilihan kantor konkret bergantung pada koordinat laporan dan keenam belas kantor pada registri \texttt{agencies\_seed.json} (Bagian~\ref{subsec:registri-instansi}), bukan jarak yang diukur pada satu lokasi tertentu.

\begin{table}[htbp]
\centering
\small
\caption{Contoh perutean instansi terdekat dari hasil klasifikasi dan koordinat laporan.}
\label{tab:demo-perutean}
\begin{tabular}{p{4.5cm} p{3.5cm} p{4.5cm}}
\toprule
\textbf{Ringkasan Laporan} & \textbf{Instansi Prediksi} & \textbf{Assignment Instansi Terdekat} \\
\midrule
Foto jalan berlubang dengan narasi ``ada lubang besar di Jalan Sudirman'' & Dinas Bina Marga & Unit Bina Marga terdekat dari koordinat laporan \\
Foto sampah menumpuk dengan narasi ``sampah dibakar di pinggir jalan'' & Satpol PP & Unit Satpol PP terdekat dari koordinat laporan \\
Foto pohon tumbang dengan narasi ``pohon roboh karena hujan deras kemarin'' & Dinas Pertamanan dan Hutan & Unit Pertamanan dan Hutan Kota terdekat dari koordinat laporan \\
\bottomrule
\end{tabular}
\end{table}

% =====================================================================
\section{Evaluasi Usabilitas Aplikasi}
\label{sec:evaluasi-usabilitas}
% =====================================================================

Selain pengujian fungsional alur pelaporan pada Bagian~\ref{sec:hasil-flutter}, aplikasi SmartCityApps dievaluasi dari sisi pengalaman pengguna untuk mengukur sejauh mana warga dapat mengoperasikan aplikasi dengan mudah, efektif, dan memuaskan. Evaluasi usabilitas penting karena nilai klasifikasi multimodal yang tinggi hanya bermanfaat apabila warga benar-benar mampu menyelesaikan alur pelaporan dari pengambilan foto hingga pengiriman laporan tanpa hambatan berarti. Dalam kerangka penelitian ini, evaluasi usabilitas merupakan bagian dari tahap \textit{Evaluasi} pada DSRM sekaligus melengkapi pemenuhan kebutuhan fungsional pada RQ1, yaitu memastikan artefak klien-server yang dibangun tidak hanya berfungsi tetapi juga dapat dioperasikan dengan nyaman oleh pengguna sasaran.

\subsection{Metode Evaluasi Usabilitas}
\label{sec:metode-usabilitas}

Usabilitas diukur menggunakan \textit{System Usability Scale} (SUS), instrumen baku yang dikembangkan oleh \citet{brookeSUSQuick1996} sebagai alat ukur usabilitas yang cepat, andal, dan dapat dibandingkan antar sistem. SUS dipilih karena tiga alasan: instrumen ini telah menjadi standar de facto pada bidang interaksi manusia--komputer, menghasilkan satu skor tunggal pada rentang 0--100 yang mudah diinterpretasikan dan di-\textit{benchmark}, serta singkat untuk diisi sehingga tidak membebani responden. SUS baku terdiri atas sepuluh pernyataan dengan skala Likert lima poin (1 = sangat tidak setuju hingga 5 = sangat setuju), yang disusun berselang-seling antara pernyataan bernada positif dan negatif untuk mengurangi bias respons. Pada penelitian ini digunakan \textbf{adaptasi sembilan butir}: lima pernyataan bernada positif (P1, P2, P4, P6, P8) dan empat pernyataan bernada negatif (P3, P5, P7, P9). Satu pernyataan negatif baku, yaitu ``saya merasa aplikasi ini terlalu rumit'', tidak disertakan pada formulir yang disebarkan; konsekuensi metodologisnya dijelaskan pada paragraf penskoran berikut dan dicatat sebagai keterbatasan pada Bab~\ref{ch:simpulan}. Agar instrumen valid bagi responden berbahasa Indonesia, kesembilan pernyataan menggunakan terjemahan Bahasa Indonesia yang mengacu pada adaptasi yang telah divalidasi oleh \citet{sharfinaIndonesianAdaptationSystem2016}.

Evaluasi melibatkan 30 responden yang memenuhi profil pengguna sasaran, yaitu Warga Pelapor yang terbiasa menggunakan \textit{smartphone} Android. Ukuran sampel ini sejalan dengan praktik evaluasi usabilitas; literatur SUS menunjukkan bahwa estimasi skor relatif stabil mulai sekitar 12--15 responden, dan tinjauan sistematis terkini menegaskan SUS tetap kokoh pada sampel kecil, sehingga sampel 30 orang memberikan estimasi yang memadai disertai interval kepercayaan yang dapat dilaporkan \citep{lewisSauroPastPresent2018, vlachogianniTseliosPerceivedUsability2022}. Setiap responden terlebih dahulu mencoba menyelesaikan serangkaian skenario tugas yang merepresentasikan alur inti pelaporan (Tabel~\ref{tab:skenario-usabilitas}) pada perangkat Android, kemudian mengisi formulir evaluasi daring (Google Form) yang memuat profil responden, pakta integritas, kesembilan pernyataan SUS (selengkapnya pada Tabel~\ref{tab:rekap-item-sus}), serta sejumlah pertanyaan umpan balik kualitatif. Rekapitulasi jawaban per responden secara anonim tersedia pada repositori penelitian (Lampiran~\ref{app:tautan}).

Skor SUS dihitung mengikuti prinsip penskoran baku \citep{brookeSUSQuick1996}. Untuk pernyataan bernada positif, kontribusi skor diperoleh dengan mengurangi 1 dari respons ($R_i - 1$); untuk pernyataan bernada negatif, kontribusi diperoleh dengan mengurangi respons dari 5 ($5 - R_j$). Setiap kontribusi karena itu bernilai 0--4. Pada SUS baku sepuluh butir, jumlah kontribusi (maksimum 40) dikalikan 2,5 agar terskala ke rentang 0--100. Karena instrumen yang digunakan hanya memuat sembilan butir, skor mentah maksimum adalah $9 \times 4 = 36$; agar skor tetap sebanding dengan \textit{benchmark} SUS pada rentang 0--100, jumlah kontribusi dinormalisasi dengan faktor $100/36 \approx 2{,}78$ menggantikan faktor 2,5 baku, sebagaimana dinyatakan pada Persamaan~\ref{eq:sus}. Dengan $R_i$ adalah respons untuk pernyataan ke-$i$, himpunan pernyataan positif adalah $\{1,2,4,6,8\}$ dan himpunan pernyataan negatif adalah $\{3,5,7,9\}$.

\begin{equation}
\label{eq:sus}
\text{Skor SUS} = \left( \sum_{i \in \{1,2,4,6,8\}} (R_i - 1) \; + \sum_{j \in \{3,5,7,9\}} (5 - R_j) \right) \times \frac{100}{36}
\end{equation}

Interpretasi skor mengacu pada tiga rujukan \textit{benchmark} yang lazim digunakan. Pertama, ambang keberterimaan \citet{bangorDeterminingWhatIndividual2009}: skor rata-rata produk berada di sekitar 68, sehingga skor di atas 68 tergolong di atas rata-rata; skor di bawah 50 tergolong \textit{tidak dapat diterima}, 50--70 \textit{marginal}, dan di atas 70 \textit{dapat diterima}. Kedua, skala kata sifat (\textit{adjective rating scale}) \citet{bangorEmpiricalEvaluation2008} yang memetakan skor ke deskriptor \textit{poor}, \textit{OK}, \textit{good}, \textit{excellent}, hingga \textit{best imaginable}. Ketiga, skala penilaian berkurva (\textit{curved grading scale}) dan peringkat persentil \citet{lewisSauroPastPresent2018}, yang menempatkan skor relatif terhadap distribusi normatif (rata-rata 68, simpangan baku 12,5); ketiga rujukan ini dikonsolidasikan dan masih dipakai pada tinjauan SUS terkini \citep{vlachogianniTseliosPerceivedUsability2022}. Konsistensi internal instrumen dilaporkan dengan koefisien \textit{Cronbach's alpha} sebagaimana praktik validasi SUS pada perangkat \textit{smartphone} di Indonesia \citep{usmanGustalikaPengujianValiditas2022}, dan presisi estimasi rata-rata dinyatakan dengan interval kepercayaan 95\% berbasis distribusi $t$ mengingat ukuran sampel yang terbatas.

Perlu ditegaskan lingkup pengukuran ini. SUS menangkap dimensi \textit{kepuasan} (\textit{satisfaction}) atau usabilitas yang dipersepsikan, bukan efektivitas dan efisiensi secara langsung; metrik objektif seperti \textit{task success rate}, \textit{time-on-task}, dan \textit{error rate} tidak diinstrumentasi terpisah pada penelitian ini meskipun setiap responden menjalankan skenario tugas pada Tabel~\ref{tab:skenario-usabilitas}. Selain itu, evaluasi terbatas pada peran Warga Pelapor; layar peran Operator Instansi dan Moderator tidak diujikan kepada pengguna (disajikan sebagai \textit{golden render} pada Bagian~\ref{sec:ui-screenshots}). Kedua batasan tersebut, beserta komposisi responden yang condong pada kelompok muda melek-digital, dicatat sebagai keterbatasan evaluasi usabilitas pada Bab~\ref{ch:simpulan}.

\begin{table}[htbp]
\centering
\small
\caption{Skenario tugas alur inti pelaporan yang dicoba setiap responden sebelum mengisi kuesioner SUS.}
\label{tab:skenario-usabilitas}
\begin{tabular}{p{1.4cm} p{10cm}}
\toprule
\textbf{Kode} & \textbf{Skenario Tugas} \\
\midrule
T-01 & Responden mendaftar akun baru atau masuk ke akun yang sudah ada. \\
T-02 & Responden membuka layar buat laporan dan mengambil foto kejadian menggunakan kamera. \\
T-03 & Responden memilih foto kejadian dari galeri perangkat. \\
T-04 & Responden mengisi narasi atau deskripsi laporan. \\
T-05 & Responden menetapkan lokasi laporan melalui GPS atau memilih titik secara manual di peta. \\
T-06 & Responden menjalankan klasifikasi AI dan meninjau hasil prediksi instansi, \textit{confidence}, serta instansi tujuan terdekat. \\
T-07 & Responden mengoreksi kategori instansi secara manual pada layar tinjauan AI. \\
T-08 & Responden mengirim laporan lalu memeriksa kemunculannya beserta status pada riwayat laporan. \\
\bottomrule
\end{tabular}
\end{table}

\subsection{Hasil Evaluasi Usabilitas}
\label{sec:hasil-usabilitas}

Evaluasi usabilitas diikuti oleh 30 responden. Profil responden dirangkum pada Tabel~\ref{tab:profil-responden}. Mayoritas responden berstatus mahasiswa (86,7\%) dengan rentang usia dominan 20--22 tahun (80,0\%, rata-rata usia 21,9 tahun) dan hampir seluruhnya tergolong pengguna \textit{smartphone} aktif (96,7\% menjawab ``sering''). Komposisi ini sesuai dengan profil pengguna sasaran aplikasi, yaitu warga melek digital yang terbiasa mengoperasikan aplikasi seluler, sekaligus menjadi konteks penting saat menafsirkan skor: responden relatif paham teknologi sehingga skor usabilitas yang diperoleh cenderung merepresentasikan batas atas pengalaman pengguna umum.

\begin{table}[htbp]
\centering
\small
\caption{Profil 30 responden evaluasi usabilitas.}
\label{tab:profil-responden}
\begin{tabular}{p{4.2cm} p{5cm} r r}
\toprule
\textbf{Karakteristik} & \textbf{Kategori} & \textbf{Jumlah} & \textbf{Persentase} \\
\midrule
Rentang usia & Di bawah 20 tahun & 2 & 6,7\% \\
             & 20--22 tahun & 24 & 80,0\% \\
             & Di atas 22 tahun & 4 & 13,3\% \\
\addlinespace
Pekerjaan / status & Mahasiswa/i & 26 & 86,7\% \\
                   & Karyawan & 3 & 10,0\% \\
                   & Wiraswasta & 1 & 3,3\% \\
\addlinespace
Frekuensi penggunaan & Sering & 29 & 96,7\% \\
\textit{smartphone}  & Kadang-kadang & 1 & 3,3\% \\
\bottomrule
\end{tabular}
\end{table}

\subsubsection{Skor SUS Keseluruhan}

Skor SUS per responden dihitung dengan Persamaan~\ref{eq:sus} lalu diringkas secara deskriptif pada Tabel~\ref{tab:sus-deskriptif}. Aplikasi SmartCityApps memperoleh \textbf{skor SUS rata-rata 71,6} (simpangan baku 14,6; median 72,2). Nilai ini berada di atas ambang rata-rata industri sebesar 68 \citep{bangorDeterminingWhatIndividual2009}, sehingga usabilitas aplikasi tergolong \textit{dapat diterima} (\textit{acceptable}) dengan deskriptor kata sifat \textit{good} \citep{bangorEmpiricalEvaluation2008} dan setara nilai \textit{B-/C+} pada skala penilaian berkurva serta kira-kira persentil ke-61 relatif terhadap distribusi normatif \citep{lewisSauroPastPresent2018}. Interval kepercayaan 95\% rata-rata adalah [66,1; 77,0]; karena batas bawah interval ini masih sedikit di bawah 68, klaim bahwa usabilitas aplikasi berada di atas rata-rata bersifat indikatif dan belum sepenuhnya terbukti secara statistik pada ukuran sampel ini. Konsistensi internal sembilan butir tergolong memadai dengan koefisien \textit{Cronbach's alpha} sebesar 0,74, sejalan dengan ambang reliabilitas yang lazim diterima ($\alpha \geq 0{,}70$) dan dengan studi reliabilitas SUS pada perangkat \textit{smartphone} di Indonesia yang melaporkan $\alpha = 0{,}835$ \citep{usmanGustalikaPengujianValiditas2022}.

\begin{table}[htbp]
\centering
\small
\caption{Statistik deskriptif skor SUS aplikasi SmartCityApps ($n = 30$).}
\label{tab:sus-deskriptif}
\begin{tabular}{l r}
\toprule
\textbf{Statistik} & \textbf{Nilai} \\
\midrule
Rata-rata (mean) & 71,6 \\
Median & 72,2 \\
Simpangan baku (SD) & 14,6 \\
Minimum & 33,3 \\
Maksimum & 97,2 \\
Interval kepercayaan 95\% (uji $t$) & [66,1; \, 77,0] \\
Cronbach's $\alpha$ (9 butir) & 0,74 \\
\midrule
Kategori keberterimaan \citep{bangorDeterminingWhatIndividual2009} & Dapat diterima \\
Deskriptor kata sifat \citep{bangorEmpiricalEvaluation2008} & \textit{Good} \\
Peringkat persentil \citep{lewisSauroPastPresent2018} & $\approx$ ke-61 \\
\bottomrule
\end{tabular}
\end{table}

Sebaran skor individual ditampilkan pada Gambar~\ref{fig:sus-distribution}. Sebanyak 17 dari 30 responden (56,7\%) memberikan skor di atas 70 (kategori \textit{dapat diterima}), 12 responden (40,0\%) berada pada rentang \textit{marginal} 50--70, dan hanya 1 responden (3,3\%) memberi skor di bawah 50. Posisi skor rata-rata terhadap ketiga skala interpretasi baku (keberterimaan, kata sifat, dan nilai huruf) divisualisasikan pada Gambar~\ref{fig:sus-interpretation}, dengan penanda rata-rata beserta rentang interval kepercayaan 95\%.

\begin{figure}[htbp]
    \centering
    \includegraphics[width=0.82\textwidth]{bab_4/sus_distribution.png}
    \caption{Distribusi skor SUS individual ($n = 30$) dengan penanda rata-rata, pita interval kepercayaan 95\%, dan garis \textit{benchmark} industri 68.}
    \label{fig:sus-distribution}
\end{figure}

\begin{figure}[htbp]
    \centering
    \includegraphics[width=\textwidth]{bab_4/sus_interpretation_scale.png}
    \caption{Posisi skor SUS rata-rata SmartCityApps (71,6) pada skala interpretasi baku: keberterimaan \citep{bangorDeterminingWhatIndividual2009}, kata sifat \citep{bangorEmpiricalEvaluation2008}, dan nilai huruf \citep{lewisSauroPastPresent2018}. Garis galat menunjukkan interval kepercayaan 95\%.}
    \label{fig:sus-interpretation}
\end{figure}

\subsubsection{Analisis Per-Pernyataan}

Untuk mengubah skor agregat menjadi temuan yang dapat ditindaklanjuti, Tabel~\ref{tab:rekap-item-sus} merinci rata-rata respons mentah (skala 1--5) beserta kontribusi terkoreksi-arah (skala 0--4, semakin tinggi semakin baik) untuk setiap pernyataan, dan Gambar~\ref{fig:sus-per-item} mengurutkannya dari yang terlemah ke terkuat. Pernyataan bernada positif memperoleh nilai konsisten tinggi: keyakinan fitur berjalan semestinya (P4, kontribusi 3,47), kemudahan penggunaan (P2, 3,33), serta keyakinan orang lain cepat memahami aplikasi (P6, 3,27) menjadi aspek terkuat. Sebaliknya, tiga kontribusi terendah seluruhnya berasal dari pernyataan bernada negatif terkait \textit{learnability}, yaitu kebutuhan bantuan orang lain atau teknisi (P3, kontribusi 1,90), kebutuhan membiasakan diri sebelum memakai aplikasi (P9, 2,20), dan persepsi ketidakkonsistenan antarmuka (P5, 2,27). Pola ini menunjukkan bahwa hambatan usabilitas utama bukan terletak pada fungsionalitas inti, melainkan pada kurva pembelajaran awal dan konsistensi tampilan, yang menjadi dasar usulan perbaikan pada Bab~\ref{ch:simpulan}.

\begin{table}[htbp]
\centering
\small
\caption{Rekapitulasi per pernyataan SUS: rata-rata respons mentah (skala 1--5) dan kontribusi skor terkoreksi-arah (skala 0--4). Tanda (+) menandai pernyataan positif dan ($-$) pernyataan negatif.}
\label{tab:rekap-item-sus}
\begin{tabular}{c p{7.6cm} c c c}
\toprule
\textbf{No.} & \textbf{Pernyataan} & \textbf{Arah} & \textbf{Mentah (1--5)} & \textbf{Kontribusi (0--4)} \\
\midrule
P1 & Akan menggunakan aplikasi ini lagi. & + & 4,27 & 3,27 \\
P2 & Aplikasi mudah digunakan. & + & 4,33 & 3,33 \\
P3 & Membutuhkan bantuan orang lain/teknisi. & $-$ & 3,10 & 1,90 \\
P4 & Fitur-fitur berjalan semestinya. & + & 4,47 & 3,47 \\
P5 & Banyak hal tidak konsisten pada aplikasi. & $-$ & 2,73 & 2,27 \\
P6 & Orang lain cepat memahami aplikasi. & + & 4,27 & 3,27 \\
P7 & Aplikasi membingungkan. & $-$ & 2,13 & 2,87 \\
P8 & Tidak ada hambatan saat menggunakan. & + & 4,20 & 3,20 \\
P9 & Perlu membiasakan diri terlebih dahulu. & $-$ & 2,80 & 2,20 \\
\bottomrule
\end{tabular}
\end{table}

\begin{figure}[htbp]
    \centering
    \includegraphics[width=0.85\textwidth]{bab_4/sus_per_item.png}
    \caption{Kontribusi skor per pernyataan SUS (skala 0--4 terkoreksi-arah), diurutkan dari pernyataan terlemah ke terkuat. Pernyataan dengan kontribusi di bawah titik tengah (2,5) menandai aspek yang paling perlu diperbaiki.}
    \label{fig:sus-per-item}
\end{figure}

\subsubsection{Analisis Sensitivitas dan Kualitas Respons}

Pemeriksaan kualitas respons menemukan tiga responden (R20, R23, R29) yang memberikan jawaban identik (nilai 5) pada seluruh pernyataan, termasuk pada pasangan pernyataan positif dan negatif yang secara logis saling bertolak belakang. Pola \textit{straight-lining} ini mengindikasikan \textit{acquiescence bias} atau pengisian yang kurang cermat, dan masing-masing menghasilkan skor 55,6. Sebagai uji ketahanan (\textit{robustness}), analisis diulang dengan mengeluarkan ketiga responden tersebut. Hasilnya justru menguatkan kesimpulan: skor rata-rata naik menjadi 73,4 (95\% CI [67,7; 79,0]) dan konsistensi internal membaik menjadi $\alpha = 0{,}81$. Dengan demikian, kesimpulan bahwa aplikasi memiliki usabilitas pada kategori \textit{dapat diterima} bersifat \textit{robust} terhadap keberadaan respons anomali, dan seluruh 30 respons tetap dipertahankan pada pelaporan utama agar tidak mengurangi keterwakilan sampel.

\subsubsection{Umpan Balik Kualitatif}

Selain penilaian kuantitatif, responden memberikan umpan balik melalui tiga pertanyaan terbuka. Pada pertanyaan kelebihan utama, mayoritas responden menyoroti kemudahan dan efisiensi proses pelaporan serta nilai tambah klasifikasi otomatis instansi (``melaporkan masalah lingkungan menjadi mudah'', ``mampu membantu proses pelaporan pada instansi terkait dengan cepat dan efisien'', ``bisa meninjau hasil prediksi instansi''). Pada pertanyaan aspek yang perlu diperbaiki, sebagian besar responden menyatakan tidak menemukan kendala berarti, sementara masukan yang spesifik mengarah pada kejelasan umpan balik (\textit{feedback}) setelah laporan dikirim. Pada pertanyaan kesediaan menggunakan aplikasi untuk melaporkan masalah lingkungan, mayoritas responden menyatakan bersedia dengan alasan kepraktisan. Rangkuman umpan balik kualitatif ini selaras dengan temuan kuantitatif, yaitu fungsionalitas inti dinilai kuat sedangkan kejelasan alur pasca-pengiriman dan kurva pembelajaran awal menjadi ruang perbaikan, dan menjadi dasar usulan penyempurnaan antarmuka pada Bab~\ref{ch:simpulan}.

% =====================================================================
\section{Pembahasan dan Analisis Kesalahan}
\label{sec:pembahasan-error}
% =====================================================================

\subsection{Pasangan Kelas yang Sering Tertukar}

\textit{Confusion matrix} ternormalisasi pada konfigurasi terbaik ditampilkan pada Gambar~\ref{fig:confusion-matrix}.

\begin{figure}[htbp]
    \centering
    \includegraphics[width=0.95\textwidth]{bab_4/confusion_matrix.png}
    \caption{\textit{Confusion matrix} ternormalisasi (per baris) pada set uji untuk konfigurasi terbaik DINOv3 large + mE5 large + CatBoost+PGS.}
    \label{fig:confusion-matrix}
\end{figure}

Analisis \textit{confusion matrix} pada konfigurasi terbaik mengidentifikasi tiga pola kesalahan yang dominan. Pertama, kelas \textbf{Kelurahan} memiliki \textit{recall} terendah di antara kelas mayoritas (0,6939), dengan laporan miliknya tersebar ke Dinas Perhubungan (134 sampel), Dinas Sumber Daya Air (110 sampel), Instansi lain (92 sampel), dan Satpol PP (86 sampel); hal ini wajar karena laporan tingkat kelurahan mencakup masalah lingkungan permukiman yang beririsan dengan banyak instansi. Kedua, kelas \textbf{Instansi lain} sebagai kategori \textit{catch-all} paling sulit dipisahkan (hanya 322 dari 663 sampel benar), karena memuat laporan heterogen yang menyerupai kelas lain. Ketiga, \textbf{Dinas Sumber Daya Air} memiliki \textit{precision} rendah (0,5577) akibat laporan banjir dan saluran air dari kelas lain (terutama Bina Marga dan Kelurahan) yang keliru diprediksi sebagai Sumber Daya Air. Berbeda dengan dugaan awal, kelas minoritas Dinas Cipta Karya justru memiliki \textit{F1} tinggi (0,8584), sehingga kesalahan lebih banyak ditentukan oleh tumpang-tindih semantik antar kelas daripada sekadar kelangkaan sampel.

\subsection{Sumber Ambiguitas}

Tiga sumber ambiguitas dapat diidentifikasi. Pertama, label noise pada tahap pemetaan 487 nama SKPD ke sembilan kategori target, yaitu beberapa laporan secara historis dapat ditugaskan ke lebih dari satu instansi tergantung kebijakan operator. Kedua, narasi yang sangat singkat tanpa kata kunci eksplisit, contohnya hanya menulis ``mohon ditangani'' tanpa menyebut jenis masalah, sehingga kontribusi modalitas teks menjadi rendah. Ketiga, citra dengan kualitas rendah seperti foto buram atau pencahayaan minim, sehingga kontribusi modalitas citra juga rendah. Pada situasi ketiga ini, kalibrasi probabilitas dengan PGS justru bermanfaat karena menampilkan ketidakpastian tinggi yang dapat memicu pengguna untuk mengoreksi kategori secara manual.

\subsection{Keterbatasan Dataset}

Selain ketidakseimbangan kelas, dataset Jakarta CRM memiliki keterbatasan lain, yaitu cakupan geografis yang hanya mencakup wilayah DKI Jakarta sehingga generalisasi ke daerah lain belum diuji, serta domain laporan yang dibatasi pada kategori publik sehingga laporan sensitif seperti yang berkaitan dengan keamanan tidak masuk ke set pelatihan. Kedua keterbatasan ini terbuka sebagai arah penelitian lanjutan, misalnya dengan memperluas dataset ke daerah lain atau menambahkan cakupan kategori yang lebih luas.

\subsection{Diskusi Keputusan Engineering}

Tiga keputusan \textit{engineering} utama pada penelitian ini menerima konfirmasi empiris dari hasil eksperimen. Pertama, keputusan untuk membekukan kedua \textit{encoder} terbukti efisien karena memungkinkan ablasi pasangan \textit{encoder} dijalankan dalam satu sesi eksperimen tanpa biaya pelatihan ulang \textit{encoder}. Kedua, keputusan memilih CatBoost sebagai \textit{classifier head} terbukti tepat karena seluruh delapan properti yang dijabarkan pada Bagian~\ref{sec:justifikasi-catboost} terealisasi pada hasil eksperimen, mulai dari biaya pelatihan rendah, kalibrasi probabilitas yang dapat diaudit, hingga paritas numerik yang terjaga pada ekspor ONNX. Inferensi \textit{classifier head} CatBoost+PGS dengan 30 ensembel virtual pada seluruh 9.266 sampel uji hanya memerlukan sekitar 0,05 ms per sampel, sehingga latensi \textit{end-to-end} server didominasi oleh tahap ekstraksi \textit{embedding} kedua \textit{encoder}, bukan oleh \textit{classifier head}. Ketiga, keputusan memilih jalur inferensi \textit{on-server} terbukti praktis karena ukuran \textit{artifact} ONNX \textit{image encoder} sebesar 1,15 GB tidak dapat dibundel ke paket aplikasi Android, sebagaimana diuraikan pada Bagian~\ref{sec:on-device-vs-on-server}.

\subsection{Latensi Inferensi dan \textit{Footprint} Sumber Daya Perangkat}
\label{sec:latensi-footprint}

Untuk menjawab konsekuensi latensi secara empiris, ketiga komponen \textit{pipeline} ONNX diukur per tahap pada satu sampel di lingkungan pengembangan berbasis CPU (rerata 20 iterasi setelah \textit{warm-up}). Hasil pada Tabel~\ref{tab:latensi} menunjukkan latensi \textit{end-to-end} sekitar 1,08 detik per sampel, dengan lebih dari 99 persen waktu dihabiskan pada kedua \textit{encoder} (\textit{image} 437 ms dan \textit{text} 639 ms) sedangkan \textit{classifier head} hanya 0,3 ms. Angka ini mengonfirmasi secara kuantitatif bahwa latensi didominasi oleh tahap ekstraksi \textit{embedding}, bukan oleh \textit{classifier}. Pada \textit{deployment} dengan akselerasi GPU, misalnya kelas RTX 4090 yang dipakai pada tahap pelatihan (Tabel~\ref{tab:runpod}), latensi kedua \textit{encoder} turun ke orde puluhan milidetik karena \textit{provider} CUDA aktif pada \texttt{serve\_model.py}. Dari sisi ketersediaan, server berfungsi sebagai titik inferensi tunggal sehingga keandalan layanan bergantung pada \textit{uptime} server; mitigasi berupa replikasi dan \textit{retry policy} diuraikan sebagai saran pada Bab~\ref{ch:simpulan}. Dari sisi biaya operasional, karena \textit{classifier head} nyaris tanpa biaya komputasi, \textit{throughput} dibatasi oleh kapasitas \textit{encoder} sehingga satu \textit{pod} GPU tunggal sudah memadai untuk beban prototipe.

\begin{table}[htbp]
\centering
\caption{Latensi inferensi per tahap pada satu sampel (diukur pada CPU lingkungan pengembangan, rerata 20 iterasi setelah \textit{warm-up}).}
\label{tab:latensi}
\begin{tabular}{l r}
\toprule
\textbf{Tahap \textit{pipeline}} & \textbf{Latensi (ms)} \\
\midrule
\textit{Image encoder} (DINOv3 \textit{large}, ONNX) & 437 \\
\textit{Text encoder} (mE5, ONNX) & 639 \\
\textit{Classifier head} (CatBoost+PGS, ONNX) & 0,3 \\
\midrule
\textbf{\textit{End-to-end}} & \textbf{1.076} \\
\bottomrule
\end{tabular}
\end{table}

Pemilihan jalur \textit{on-server} juga menjaga \textit{footprint} sumber daya pada perangkat pengguna tetap minimum, sesuai tuntutan rumusan masalah pertama. Paket aplikasi Android tidak memuat model apa pun (direktori \texttt{assets/models} kosong); berkas APK \textit{debug} berukuran 216 MB (versi rilis lebih kecil setelah \textit{tree-shaking} dan \textit{split-per-ABI}), jauh di bawah total sekitar 7,3 GB artefak ONNX (\textit{image encoder} DINOv3 \textit{large} 1,1 GB, \textit{text encoder multilingual-E5 large} 6,2 GB, dan \textit{classifier head} CatBoost 31 MB) yang harus disertakan apabila inferensi dilakukan sepenuhnya \textit{on-device}. Pada saat operasi, perangkat hanya mengirim satu permintaan \textit{multipart} berisi citra terkompresi berukuran ratusan kilobyte beserta narasi teks singkat, kemudian menerima respons JSON ringan; tidak ada komputasi model di sisi perangkat sehingga konsumsi memori, penyimpanan, dan daya pada perangkat pengguna tetap rendah.
